\documentclass[12pt]{article}
\usepackage[utf8]{inputenc}
\usepackage[T1]{fontenc}
\usepackage{hyperref}
\title{Mbo Language Resources for K-12 Curriculum Development}
\author{Compiled by Claude Sonnet 5.5 (claude-sonnet-5-5) for Andr\'as Kornai}
\date{}
\begin{document}
\maketitle

\section{Introduction}
Mbo (also Imbo or Kimbo; ISO 639-3 zmw; Glottocode mbod1238) is a Bantu language (zone D.30) of the Democratic Republic of the Congo. It is spoken in the Ituri forest area, in Mambasa Territory. A figure of about 11,000 speakers dates from 1994, and the Mbo people number about 17,000. It should not be confused with Mbo of Zambia or with Mbochi (Mbosi) of the Republic of the Congo.

\section{Primary Educational Resources}
No primers, textbooks or school materials were found in the searches.

\section{Linguistic and Reference Materials}
\begin{itemize}
\item Wikipedia article on the language: \url{https://en.wikipedia.org/wiki/Mbo_language_(Congo)}
\item Wikipedia article on the Mbo people: \url{https://en.wikipedia.org/wiki/Mbo_people_(Congo)}
\item ASJP comparative word list for Mbo (DRC). This is a short comparative list, not a dictionary: \url{https://asjp.clld.org/languages/MBO_DRC}
\item Glottolog entry: \url{https://glottolog.org/resource/languoid/id/mbod1238.bigmap.html}
\end{itemize}

\section{Religious and Community Resources}
None found.

\section{Digital and Audio Resources}
None found. A language-code page exists at \url{https://localizely.com/language-code/zmw}, but it has no learning content.

\section{Curriculum Development Context}
The language is vigorous in its community, but the searches found no teaching materials and no school programme. No Mbo dictionary or grammar turned up. A curriculum would therefore need new fieldwork and an orthography, developed with Mbo speakers. Working with local education authorities and SIL or Bantu linguists could help.
\end{document}